\section{What remains after effective finiteness}
\label{sec:outlook}

The results separate two questions that are often conflated. At each
fixed occurrence count, equality is completely decidable by a finite
calculation with an explicit polynomial bound. To settle Markov
uniqueness, one must control the entire sequence of these finite
questions as the count grows. Theorem~\ref{thm:finite-decision} answers
the first question and supplies a precise setting for the second.

The restriction is geometric, not just computational. The residue orbit
\eqref{eq:residue-orbit} is invariant under a common change of marking
of the pair. A change of basis can simplify the matrix character or the
word, but cannot in general reduce the least representative of that
orbit. Thus the method does not become a proof of uniqueness by
repeatedly choosing a new convenient basis. The remaining intersection
range is also unbounded: the result through $24$ is a uniform theorem
at those intersections, not an estimate on all possible intersections.

\begin{example}[Isometric pairs at every occurrence count]
\label{ex:unbounded-isometric-family}
For any $h\geqslant2$, take root homology columns
$v=(h+1,1)$ and $b=(-1,0)$. They form a simple marking since
$\det(v,b)=1$. The actual order-three action $G$ from
Proposition~\ref{prop:isometry-group} satisfies
$$
 Gv=(-1,h)=hv+(h^2+h+1)b.
$$
Thus an isometric equal-length pair has original type
$(k,t)=(h+1,1)$ at every occurrence count. For
$N=h^2+h+1$, the sign-and-inversion orbit is
$\{h,h+1,N-h,N-h-1\}$, whose least positive representative is $h$.
These pairs therefore retain unbounded least occurrence count under
common changes of marking.
\end{example}

\subsection{An unbounded family beyond the finite cover}

Proposition~\ref{prop:intersection-twentythree} is the second member of a
family, and the family is infinite. Keep $X=AB^3$ and $Y=AB^4$, and for
$n\geqslant0$ compare $A$ with $X(X^nY)^2$. The pair $(X,X^nY)$ is again a
simple basis, so the second word is primitive and simple; its homology is
$(2n+3,6n+11)$, and its two long gaps are separated by $n$ and $n+1$ short
ones, so the cyclic gap pattern is mechanical. Both the occurrence count and
the intersection number therefore grow without bound along the family. At
$n=0$ the pair is the three-occurrence type $(k,t)=(3,2)$, settled by
\cref{thm:small-occurrences}; at $n=1$ it is the residue-five pair at
intersection seventeen, also settled there; at $n=2$ it is
\cref{prop:intersection-twentythree}.

We have proved that $\tr A\ne\tr(X(X^nY)^2)$ for every $n\geqslant0$ and every
positive integral marked character, so that the residue $h$ at intersection
$3h+2$ carries no equal-length pair for any odd $h\geqslant3$. The argument
beyond $n=2$ is not the one used here: it replaces the fixed elimination by a
spectral representation of the fixed-coordinate ray, with uniform bounds on
its terminal core, and it is accompanied by an elementary congruence filter
modulo nine which by itself already removes every quotient divisible by four
and every quotient congruent to six modulo twelve. Because that machinery is
of a different kind and about as long as the present article, it is written up
separately. We record the statement here because it shows what the finite
cover of \cref{thm:small-intersections} does not see: an unbounded prescribed
family can be excluded outright, while the cover itself stops at twenty-four.

This does not bound the occurrence count of an arbitrary equality, and it does
not cover every residue at the intersections it meets.

The error estimates point to a more specific analytic question. An
isometric equality can occur inside the finite domain and must be
retained. Any uniform refinement must therefore distinguish the actual
symmetry equations from a possible nonisometric equality. Excluding all
equalities at large occurrence count would be too strong. A promising
next step is to characterize equality in the dominant-state and
Euclidean-descent constraints in terms of the twelve actions of
Proposition~\ref{prop:isometry-group}. The present bounds establish
neither that characterization nor a bound on the first possible
nonisometric occurrence count.

There is a second reformulation of the same uniqueness question, arithmetic
rather than geometric, and it is worth recording here because it shows from a
different side why a bound on the occurrence count is the missing ingredient
and not a technicality. A simple curve is presented by a Christoffel word, and
Frobenius attached to the corresponding Markov number $m$ a square root of $-1$
modulo $m$~\cite{Frobenius1913}. That root is not an abstract residue: it is a
letter count. Writing a proper Christoffel word as $a\operatorname{Pal}(d)b$
for a directive word $d$, with $q$ letters $a$ and $s$ letters $b$, the word
directed by the doubled directive $\widehat d\,d$, where $\widehat d$ is the
reversal complement, has length exactly $m=q^2+s^2$ and precisely $sq^{-1}$
letters $b$; and, the incidence matrix of a doubled directive being a Gram
matrix and hence symmetric, that same number says where the word splits in its
standard factorization into two Christoffel words. The classical dictionary
behind these statements is the involution theory of Christoffel words and
Sturmian morphisms~\cite[Section~3]{BertheDeLucaReutenauer2008} together with
Reutenauer's correspondence~\cite{Reutenauer2009}; the proofs of the two
statements just quoted are given in a companion article of this series. A
combinatorial proof that a Markov number is a sum of two squares was obtained
independently by L.~Vuillon and A.~De~Luca (personal communication, 2026).

Read through that dictionary, uniqueness becomes the assertion that a Markov
number determines the primitive representation $m=q^2+s^2$ of every occurrence
carrying it. Since $m$ has $2^{\omega(m)-1}$ primitive representations, with
$\omega$ the number of distinct prime divisors, the question is empty when $m$
is a prime power or twice one, which recovers the classical results in that
range and explains why they stop there. What it does not supply is any control
as $\omega(m)$ grows, and in particular no bound on the occurrence count. The
reformulation and the present theorems therefore fail in the same place, from
opposite directions: the arithmetic side decides a single label and cannot
follow a family, while \cref{thm:finite-decision} decides a whole family at a
fixed count and cannot follow the count.

The contribution of the method is that the trace comparison keeps
three kinds of structure visible at once: the positive paths of the
word, the integral Vieta orbit of its character, and the marking that
recognizes a surface symmetry. The positive paths and integral Vieta
orbit produce the terminal contradiction and hence the parameter bounds;
the marking identifies the symmetries among the retained equalities.
Together they reduce an unbounded search in matrix entries and gap
lengths to an explicit family of finite geometric decisions, while
leaving the remaining uniform question identifiable.
